% Generated draft: requires mathematical and source review.
\documentclass[11pt]{article}
\usepackage[margin=25mm]{geometry}
\usepackage{fontspec}
\setmainfont{Latin Modern Roman}
\usepackage{amsmath,amssymb,mathtools,unicode-math}
\setmathfont{Latin Modern Math}
\usepackage{enumitem,needspace,xurl,hyperref}
\setlength{\parindent}{0pt}
\setlength{\parskip}{5pt}
\setlength{\emergencystretch}{3em}
\newcommand{\N}{\mathbb{N}}
\newcommand{\Z}{\mathbb{Z}}
\newcommand{\Q}{\mathbb{Q}}
\newcommand{\R}{\mathbb{R}}
\newcommand{\C}{\mathbb{C}}
\newcommand{\F}{\mathbb{F}}
\newcommand{\A}{\mathbb{A}}
\newcommand{\PP}{\mathbb{P}}
\newcommand{\E}{\mathbb{E}}
\newcommand{\eps}{\varepsilon}
\newcommand{\dd}{\,\mathrm{d}}
\begin{document}
\section*{UnsolvedMath ID: 30005549}
\textbf{Source problem code:} catalogue entry
\subsection{Definitions and mathematical statement}
Let \(M\subset\R^n\) be a compact \(C^1\) hypersurface, as in the supplied record. For \(Q\geq1\) and \(\delta>0\), define
\[
N_M(\delta,Q)=\#\left\{(q,p)\in\Z\times\Z^n: Q\leq q<2Q,\quad \operatorname{dist}(p/q,M)\leq\frac{\delta}{q}\right\},
\]
where \(\operatorname{dist}(x,M)=\inf_{y\in M}\|x-y\|_2\). Thus denominators are positive integers, the denominator interval is half-open, and the distance bound includes equality. Pairs are counted without a coprimality condition; distinct representations of the same rational point are counted separately.

The supplied problem asks for a characterization of the local obstructions to the uniform estimate
\[
N_M(\delta,Q)\asymp\delta Q^n,
\qquad Q\geq1,\qquad Q^{\varepsilon-1}<\delta<\tfrac12.
\]
Under the conventional interpretation that \(\varepsilon>0\) is fixed and the constants may depend on \(M\) and \(\varepsilon\), this estimate means that there exist constants \(0<c_{M,\varepsilon}\leq C_{M,\varepsilon}<\infty\) such that
\[
c_{M,\varepsilon}\delta Q^n\leq N_M(\delta,Q)\leq C_{M,\varepsilon}\delta Q^n
\]
for every admissible \(Q\) and \(\delta\). This is an explicit interpretation requiring review: the supplied record does not specify the quantification or range of \(\varepsilon\), the permitted dependence of the comparison constants, or whether \(Q\) ranges over real scales or integers. The interval for \(\delta\) may be empty, in which case it imposes no condition.

The intended characterization should identify which local features prevent these bounds and under what hypotheses their absence guarantees the bounds. The record does not define a local obstruction or supply a necessary-and-sufficient criterion. It also leaves the admissible ambient dimensions and conventions concerning nonemptiness and boundary unspecified. These qualifications remain unresolved rather than being added as hypotheses.

The supplied original excerpt discusses nonvanishing Gaussian curvature, expressed through a graph Hessian, and flat or rough hypersurfaces with isolated vanishing or blow-up of Hessian determinants and other derivatives. It refers to an additional counting term for such singular behavior. These remarks concern more specialized settings: a Hessian is not supplied by the stated \(C^1\) regularity alone, and the excerpt does not provide the graph-counting definition, equation (3), Theorem 3, or the additional term. Consequently, their precise relationship to the displayed compact-hypersurface formulation cannot be reconstructed from the supplied evidence.
\subsection{Short English statement}
Count integer pairs representing rational points within distance \(\delta/q\) of a compact \(C^1\) hypersurface, with \(Q\leq q<2Q\). The task is to characterize the local obstructions to having this count uniformly comparable to \(\delta Q^n\) when \(Q^{\varepsilon-1}<\delta<1/2\). The parameter quantifiers, dimensional and boundary conventions, and connection to the truncated source excerpt require clarification. This is a statement-only draft requiring human mathematical review.
\subsection{Sources}
\begin{itemize}
\item Supplied local UnsolvedMath registry record, permanent ID \(30005549\), title \emph{Rational Point Counts near Hypersurfaces}, source problem code \texttt{catalogue entry}, in the Oberwolfach Reports Open Problems collection. The record gives proposed year 2023 and supplies no named proposer. Its labels \texttt{partially\_solved} and \texttt{corrected\_verified}, including its assertion of verification against the official report, are supplied registry metadata, not independent verification here.
\item \emph{Incidence Problems in Harmonic Analysis, Geometric Measure Theory, and Ergodic Theory} (2023), Oberwolfach Reports, pp.\ 1397--1452, DOI 10.4171/owr/2023/25: \url{https://doi.org/10.4171/owr/2023/25}. The supplied original excerpt contains a page marker for p.\ 1435 and is incomplete.
\item Supplied literature source labeled ``OWR 2023 25'': \url{https://publications.mfo.de/bitstream/handle/mfo/4055/OWR_2023_25.pdf?sequence=1}.
\item Supplied literature source labeled ``Srivastava Technau'': \url{https://arxiv.org/abs/2305.01047}. The record supplies no fuller bibliographic identification.
\item Literature assessment checked at \texttt{2026-08-22}: supplied dated curation, not independently verified. It labels the problem partially solved and reports sharp estimates for curved hypersurfaces and broad homogeneous classes, while arbitrary local obstructions and the general expected count remain unresolved. The background explicitly describes this as dated literature triage derived from the supplied OWR curation and requests independent verification of time-sensitive status and cited claims.
\end{itemize}
\end{document}
